\section{The combination problem}
\label{sec:combination}

Three axes with two settings each give eight configurations, and no deployed network is
a corner. A roster with an introduction protocol is public on formation for the agents
it can reach through a broker and private for everyone else; an enterprise workspace is
private on formation and fully open on semantics. The design question is not which
regime wins but \textbf{which configuration a given setting should hold}, and that
requires each axis to be priced separately---which is what \S\ref{sec:quantifying}
buys.

\paragraph{The axes are not equally live.} A configuration search is only interesting if
more than one dimension moves the outcome. Our measurements so far
(\S\ref{sec:evidence}) put the three weights at roughly $0.23$, $0.08$ and $0$ on
requirement $F_1$, which collapses the search: pick a point on formation and the other
two are free. We report that collapse rather than dressing it as an optimisation,
because it is the honest reading of one task family---and because it makes the next
question sharp.

\paragraph{Scenarios are the test of whether the weights are constants.} The collapse
above was measured on a single task class: facts held by different agents, aggregated
into one deliverable. That class is maximally favourable to formation---reach is
literally what it asks for---and maximally unfavourable to semantics, since there is no
shared artifact to co-edit. If the three weights are properties of the framework, they
should hold elsewhere. If they are properties of the task, they will move, and the
combination problem becomes real.

\begin{table}[t]
\centering\small
\caption{Three settings, and what each predicts for the weight on each axis. These are
predictions the framework makes and does not yet have evidence for; the task families
are specified in Appendix~\ref{app:scenarios}.}
\label{tab:scenarios}
\begin{tabular}{@{}p{0.20\linewidth}p{0.30\linewidth}p{0.13\linewidth}p{0.13\linewidth}p{0.13\linewidth}@{}}
\toprule
\textbf{Setting} & \textbf{What a task looks like} & \textbf{formation} & \textbf{semantics} & \textbf{attributes} \\
\midrule
Enterprise team & Co-edit one artifact; the people are known; the context is the
  bottleneck & low & \textbf{high} & low \\
\addlinespace[2pt]
Personal errand & Find the one counterpart who can do a specific thing; no
  relationship needed after & \textbf{high} & low & medium \\
\addlinespace[2pt]
Personal social & Reach a person through someone; the introduction is the product &
  \textbf{high, brokered} & low & \textbf{high} \\
\bottomrule
\end{tabular}
\end{table}

The enterprise row is the sharpest test, because it predicts the reverse of what we
measured: when everyone needed is already reachable, formation has no work left to do,
and the shared store that cost 11k tokens for a small gain in a pooling task should pay
for itself when the deliverable is edited jointly. A framework that only ever produces
``formation dominates'' is a restatement of the task; one that predicts where formation
stops dominating, and is then held to it, is a framework.

\paragraph{What an optimum means here.} An optimum is a configuration, not a size. On
formation it is a \emph{funnel depth}: how far past the roster a requester should be
allowed to reach---own roster, one brokered hop, two hops, open directory---given that
each additional layer widens reach and weakens recourse. The marginal return on a layer
is the number of \emph{independent} sources it adds, which saturates; the marginal cost
is what a claim from that layer is worth, which decays. The crossing is the depth. Stated
this way the optimum is measurable per scenario rather than argued per regime.
